
\chapter*{Apêndices}
\addcontentsline{toc}{chapter}{Apêndices}

\section*{A. Plano de 100 dias (pequenos passos, grande fidelidade)}
\subsection*{Como usar}
Reserve \textbf{15 minutos/dia} (ou 30 em família). Se cair, retome do dia seguinte. Cada semana tem um \textbf{tema} e metas simples. Ao final dos 100 dias, você terá consolidado hábitos que permanecem.

\subsection*{Semanas e metas}
\begin{enumerate}[leftmargin=*,nosep]
  \item \textbf{Semanas 1--2: Atenção e Palavra} — Primeiro olhar do dia para a Bíblia; Salmos na mesa (cap.~13).
  \item \textbf{Semanas 3--4: Higiene informacional} — Teste 4P; checklist de 60s; dois blocos de rede/dia (cap.~3 e 6).
  \item \textbf{Semanas 5--6: Regra de vida digital \& família} — refeição sem tela; liturgia doméstica 15 min (cap.~13).
  \item \textbf{Semanas 7--8: Identidade \& segurança digital} — mapa de contas, 2FA, plano B offline (cap.~7).
  \item \textbf{Semanas 9--10: Missão digital \& serviço} — vídeo 3T/semana; responder DMs; visita/serviço (cap.~16).
  \item \textbf{Semanas 11--12: Resiliência doméstica} — água e alimentos 7 dias; auditoria do bairro (cap.~10).
  \item \textbf{Semanas 13--14: Generosidade \& ética no trabalho} — liturgia de orçamento; linhas vermelhas (cap.~15 e 17).
  \item \textbf{Dias 99--100: Revisão e compromisso} — revisar aprendizados; carta à família (1 página) com protocolos e promessas.
\end{enumerate}

\subsection*{Roteiro diário (exemplo simples)}
\begin{checklist}
\begin{itemize}[leftmargin=*,noitemsep]
  \item \textbf{5 min Palavra}: salmo em voz alta; uma frase para lembrar.
  \item \textbf{5 min passo prático}: tarefa da semana (ex.: ativar 2FA; montar kit de acesso).
  \item \textbf{5 min oração}: gratidão (3 itens) + pedido (1 item) + Pai-Nosso.
\end{itemize}
\end{checklist}

\subsection*{Marcos semanais (checklist)}
\begin{checklist}
\begin{itemize}[leftmargin=*,noitemsep]
  \item \textbf{Semana 1}: desinstalar 1 app que só distrai.
  \item \textbf{Semana 2}: primeira refeição diária sem tela.
  \item \textbf{Semana 3}: ativar 2FA nos e-mails.
  \item \textbf{Semana 4}: ensaiar checklist 60s com a família.
  \item \textbf{Semana 5}: organizar documentos e contatos no papel.
  \item \textbf{Semana 6}: liturgia doméstica 2x na semana.
  \item \textbf{Semana 7}: montar kit de acesso (backups de recuperação).
  \item \textbf{Semana 8}: auditoria de permissões de apps.
  \item \textbf{Semana 9}: gravar 1 vídeo 3T de 60s.
  \item \textbf{Semana 10}: responder todas as DMs em 24h.
  \item \textbf{Semana 11}: água e alimentos para 7 dias.
  \item \textbf{Semana 12}: auditoria do bairro (3 mercados/2 farmácias).
  \item \textbf{Semana 13}: ajustar orçamento (três caixas) e margem de misericórdia.
  \item \textbf{Semana 14}: escrever carta de protocolos (família/igreja).
\end{itemize}
\end{checklist}

\bigskip
\section*{B. Glossário tecno-bíblico (conciso)}
\begin{description}
  \item[Anticristo] Pessoalidade e sistema de oposição a Cristo (2~Ts~2; Ap~13).
  \item[Apocalipse] ``Revelação'' de Jesus Cristo; não apenas catástrofes (Ap~1:1).
  \item[BCI] Interface cérebro-dispositivo para comunicação/controle motor.
  \item[CBDC] Moeda digital de banco central; pode ter regras de uso programáveis.
  \item[CRISPR] Ferramenta de edição genética com usos terapêuticos específicos.
  \item[Deepfake] Mídia sintética hiper-realista gerada por IA.
  \item[Drex] Arquitetura brasileira de liquidação/ativos tokenizados com dinheiro do BC na camada de atacado.
  \item[Engano digital] Produção/curadoria de narrativas falsas amplificadas por algoritmos.
  \item[Escatologia] Doutrina das últimas coisas: volta de Cristo, juízo, nova criação.
  \item[Falso Profeta] Agente de propaganda/milagres a serviço da Besta (Ap~13).
  \item[Framework NIC] Lente \textbf{N}arrativa, \textbf{I}dentidade, \textbf{C}ontrole para discernir textos e tecnologias.
  \item[IoT] Internet das Coisas; sensores/redes em objetos e infraestrutura.
  \item[JAM-DEX / eNaira / e-CNY] Exemplos de CBDCs lançadas/piloto (Jamaica, Nigéria, China).
  \item[KYC] ``Conheça seu cliente'' — procedimentos de verificação de identidade.
  \item[Marca da besta] Sinal de \textbf{lealdade} que habilita \emph{compra e venda} (Ap~13:16--17).
  \item[Milênio] Período de Ap~20; interpretado de formas diversas (ami/pre/pós).
  \item[mBridge] Plataforma multi-BCs para pagamentos transfronteiriços com CBDCs.
  \item[Passkey] Autenticação baseada em chaves criptográficas e biometria do dispositivo.
  \item[Preterismo/Futurismo/Idealismo/Historicismo] Quatro lentes de leitura de Apocalipse (cap.~18).
  \item[Ransomware] Sequestro de sistemas/dados com pedido de resgate.
  \item[Selo do Espírito] Confirma filiação em Cristo; identidade recebida (Ef~1:13).
  \item[Smart city] Cidade com coleta/uso intensivo de dados para gestão/serviços.
  \item[Transumanismo] Movimento/ideias que defendem ampliar capacidades humanas por tecnologia.
  \item[Vigilância] Monitoramento de pessoas/dados; pode servir segurança ou abuso.
\end{description}

\bigskip
\section*{C. Guia de referências bíblicas (por tema)}
\subsection*{Cristo no centro}
Ap 1; Ap 5; Ap 19:11--16; Cl 1:15--20; Hb 1.

\subsection*{Verdade e engano}
1~Jo 4:1--6; Mt 24:4--25; 2~Ts 2:9--12; Ef 4:14--15; Pv 18:17.

\subsection*{Esperança e perseverança}
Rm 5:3--5; Tg 1:2--4; Hb 6:19; 1~Pe 1:3--9; Ap 21:1--5.

\subsection*{Economia e justiça}
Pv 11:1; Mq 6:11; Tg 5:1--6; At 2:42--47; 2~Co 8--9; Ap 18.

\subsection*{Identidade e selo}
Ef 1:13--14; 2~Tm 2:19; Ap 7; Ap 14:1.

\subsection*{Missão e testemunho}
Mt 28:18--20; At 1:8; 1~Pe 3:15; 2~Ts 3:1.

\subsection*{Santidade e sobriedade}
1~Pe 5:8--10; 1~Ts 5:4--11; 2~Pe 3:11--12; Rm 12:1--2.

\subsection*{Família e comunidade}
Ef 5--6; Cl 3:12--17; Hb 10:24--25; At 20:20.

\subsection*{O Anticristo e a besta}
Dn 7; 2~Ts 2; Ap 13; Ap 17--19.

\subsection*{Nova criação}
Is 65:17--25; Rm 8:18--25; Ap 21--22.
